\documentclass[12pt,a4paper]{article}
\usepackage[margin=1in]{geometry}
\usepackage{graphicx}
\usepackage{amsmath,amssymb}

% ======================================================================
\begin{document}

% -- Title Page --------------------------------------------------------
\begin{titlepage}
\centering
\vspace*{2cm}
{\Huge\bfseries Phone-Based Inertial Navigation Unit\\[0.3cm]
Using Digital Signal Processing\par}
\vspace{1cm}
{\Large A ROS\,2 Implementation with Butterworth Filtering,\\
Madgwick Sensor Fusion, and Real-Time Web Dashboard\par}
\vspace{2cm}
{\large\textbf{Course:} Digital Signal Processing\par}
\vspace{0.5cm}
{\large\textbf{Submitted by:}\par}
\vspace{0.3cm}
{\Large Vaibhav (23f3000074)\par}
\vspace{0.3cm}
{\large BS Electronic Systems\par}
{\large Indian Institute of Technology Madras\par}
\vspace{1.5cm}
{\large\textbf{Date:} May 2026\par}
\vfill
{\small Email: \texttt{23f3000074@ds.study.iitm.ac.in}\\
Repository: \texttt{https://github.com/Bishu-crypto/phone-imu-bridge}}
\end{titlepage}

% -- Abstract ----------------------------------------------------------
\begin{abstract}
This project implements a real-time \textbf{Inertial Navigation Unit (INU)} that transforms a smartphone into an inertial measurement unit (IMU) using its built-in accelerometer and gyroscope. The system streams 6-axis IMU data via HTTP to a ROS\,2 processing pipeline that applies: (1)~\textbf{Madgwick gradient-descent sensor fusion} for orientation estimation, (2)~a \textbf{2nd-order Butterworth IIR low-pass filter} designed via the bilinear transform for noise suppression, (3)~quaternion-based \textbf{gravity removal}, (4)~\textbf{trapezoidal numerical integration} for velocity and displacement estimation, and (5)~\textbf{Zero-Velocity Updates (ZUPT)} for drift mitigation. A companion web dashboard provides real-time visualization of orientation, acceleration, navigation state, and FFT-based spectral analysis. All DSP algorithms are implemented from first principles without reliance on \texttt{scipy.signal}.
\end{abstract}

\tableofcontents
\newpage

% ======================================================================
\section{Introduction}

\subsection{Motivation}
Modern smartphones contain MEMS (Micro-Electro-Mechanical Systems) inertial sensors---accelerometers and gyroscopes---that rival dedicated low-cost IMU hardware. This project leverages these sensors to build a complete inertial navigation pipeline, demonstrating core DSP concepts including digital filter design, sensor fusion, spectral analysis, and numerical integration.

\subsection{Objectives}
\begin{enumerate}
  \item Stream raw 6-axis IMU data from an Android phone to a Linux workstation via HTTP.
  \item Apply Madgwick AHRS sensor fusion to estimate 3D orientation as a quaternion.
  \item Design and implement a 2nd-order Butterworth IIR low-pass filter from first principles using the bilinear transform.
  \item Remove gravitational acceleration using quaternion-derived rotation matrices.
  \item Estimate velocity and displacement via trapezoidal integration with ZUPT drift correction.
  \item Perform real-time FFT-based spectral analysis for noise characterization.
  \item Visualize all outputs on a real-time web dashboard.
\end{enumerate}

\subsection{Technology Stack}
\begin{itemize}
  \item \textbf{Platform:} Ubuntu 22.04, ROS\,2 Humble
  \item \textbf{Language:} Python 3.10, JavaScript (ES6)
  \item \textbf{Libraries:} NumPy, \texttt{websockets}
  \item \textbf{Phone App:} Android Sensor Logger with HTTP Push
  \item \textbf{Dashboard:} HTML5 Canvas, WebSocket client
\end{itemize}

% ======================================================================
\section{System Architecture}

The system follows a modular, publish-subscribe architecture built on ROS\,2. Five independent nodes form a DSP processing pipeline:

\begin{verbatim}
  +----------+       +---------+       +-----------+       +----------+
  |   HTTP   | raw   |Madgwick | fused | Inertial  | vel/  | WebSocket|
  | Receiver |------>| Filter  |------>|Navigation |------>|  Bridge  |
  | (phone)  |       | (AHRS)  |       |(LPF+ZUPT) | pos   |(dashboard)
  +----------+       +---------+       +-----------+       +----------+
       |                                                        ^
       |              +-----------+                             |
       +------------->| Spectral  |-----------------------------+
                      | Analysis  |  PSD (JSON)
                      +-----------+
\end{verbatim}

\bigskip
\noindent\textbf{ROS\,2 Topics:}
\begin{center}
\begin{tabular}{|l|l|l|}
\hline
\textbf{Topic} & \textbf{Type} & \textbf{Description} \\
\hline
\texttt{/phone/imu/data\_raw}   & \texttt{sensor\_msgs/Imu}    & Raw 6-axis IMU \\
\hline
\texttt{/phone/imu/data\_fused} & \texttt{sensor\_msgs/Imu}    & Fused orientation \\
\hline
\texttt{/phone/nav/velocity}    & \texttt{Vector3Stamped}      & Estimated velocity \\
\hline
\texttt{/phone/nav/position}    & \texttt{Vector3Stamped}      & Estimated displacement \\
\hline
\texttt{/phone/imu/psd}         & \texttt{std\_msgs/String}    & PSD as JSON \\
\hline
\end{tabular}
\end{center}

% ======================================================================
\section{DSP Theory and Implementation}

\subsection{Butterworth Low-Pass Filter}

\subsubsection{Why Butterworth?}
The Butterworth filter provides a \textbf{maximally flat magnitude response} in the passband---no ripple---making it ideal for IMU data where smooth attenuation of high-frequency MEMS noise is needed without passband distortion.

\subsubsection{Analog Prototype}
The 2nd-order Butterworth analog transfer function is:
\begin{equation}
H_a(s) = \frac{\omega_c^2}{s^2 + \sqrt{2}\,\omega_c\,s + \omega_c^2}
\end{equation}
This yields a $-40$\,dB/decade rolloff beyond $\omega_c$.

\subsubsection{Bilinear Transform}
To convert from continuous-time to discrete-time:
\begin{equation}
s = \frac{2}{T}\cdot\frac{z-1}{z+1}
\end{equation}
where $T = 1/f_s$ is the sampling period.

\subsubsection{Frequency Pre-Warping}
The bilinear transform introduces frequency warping. To place the $-3$\,dB point at exactly $f_c$:
\begin{equation}
\omega_c' = \frac{2}{T}\tan\!\left(\frac{\pi f_c}{f_s}\right)
\end{equation}

\subsubsection{Discrete Transfer Function}
With $K = 2f_s$:
\begin{align}
D &= K^2 + \sqrt{2}\,\omega_c'\,K + {\omega_c'}^2 \\
b_0 &= {\omega_c'}^2 / D,\quad b_1 = 2\,{\omega_c'}^2 / D,\quad b_2 = {\omega_c'}^2 / D \\
a_1 &= (2\,{\omega_c'}^2 - 2K^2) / D,\quad a_2 = (K^2 - \sqrt{2}\,\omega_c'\,K + {\omega_c'}^2) / D
\end{align}

\subsubsection{Difference Equation (Direct Form I)}
\begin{equation}
y[n] = b_0\,x[n] + b_1\,x[n\!-\!1] + b_2\,x[n\!-\!2] - a_1\,y[n\!-\!1] - a_2\,y[n\!-\!2]
\end{equation}
With $f_c = 5$\,Hz and $f_s = 100$\,Hz, this removes sensor noise while preserving human-scale motion.

\subsubsection{IIR vs FIR Trade-off}
IIR was chosen over FIR because: (1)~lower order means fewer coefficients and less delay, (2)~non-linear phase is acceptable for IMU navigation where exact phase alignment is not critical.

\subsubsection{Implementation (dsp\_filters.py)}
\begin{verbatim}
def _design(self):
    fs, fc = self.sample_rate, self.cutoff_hz
    wc = 2.0 * fs * math.tan(math.pi * fc / fs)  # Pre-warp
    k  = 2.0 * fs
    k2, wc2 = k * k, wc * wc
    sqrt2_wc_k = math.sqrt(2.0) * wc * k
    denom = k2 + sqrt2_wc_k + wc2
    b0 = wc2 / denom;  b1 = 2.0 * wc2 / denom;  b2 = b0
    a1 = (2.0*wc2 - 2.0*k2) / denom
    a2 = (k2 - sqrt2_wc_k + wc2) / denom
    return (b0, b1, b2), (1.0, a1, a2)
\end{verbatim}

% ----------------------------------------------------------------------
\subsection{Madgwick AHRS Sensor Fusion}

\subsubsection{Problem Statement}
Given noisy accelerometer ($\mathbf{a}$) and gyroscope ($\boldsymbol{\omega}$) readings, estimate device orientation as a unit quaternion $\mathbf{q} = [w, x, y, z]^T$.

\subsubsection{Algorithm}
The Madgwick filter uses \textbf{gradient descent} to correct gyroscope-integrated orientation toward the gravity reference:
\begin{equation}
\mathbf{q}_{t+1} = \mathbf{q}_t + \dot{\mathbf{q}}\,\Delta t, \qquad
\dot{\mathbf{q}} = \tfrac{1}{2}\mathbf{q}_t \otimes \boldsymbol{\omega} - \beta\,\frac{\nabla f}{\|\nabla f\|}
\end{equation}

\subsubsection{Objective Function}
\begin{align}
f_1 &= 2(q_1 q_3 - q_0 q_2) - a_x \\
f_2 &= 2(q_0 q_1 + q_2 q_3) - a_y \\
f_3 &= 2(0.5 - q_1^2 - q_2^2) - a_z
\end{align}
The Jacobian $\mathbf{J}$ is computed analytically; gradient is $\nabla f = \mathbf{J}^T \mathbf{f}$.

\subsubsection{Gain Parameter $\beta$}
$\beta = 0.033$ balances gyroscope trust (low $\beta$: fast response) vs.\ accelerometer trust (high $\beta$: drift-free but noisy).

\subsubsection{Startup Calibration}
First 50 samples are averaged to compute an initial quaternion from the gravity vector via pitch/roll decomposition.

% ----------------------------------------------------------------------
\subsection{Gravity Removal}

The accelerometer measures \textbf{specific force}:
\begin{equation}
\mathbf{a}_{\text{meas}} = \mathbf{a}_{\text{linear}} + \mathbf{R}^T \mathbf{g}
\end{equation}
We remove gravity using the quaternion-derived rotation matrix:
\begin{align}
\mathbf{a}_{\text{linear,body}} &= \mathbf{a}_{\text{meas}} - \mathbf{R}^T \mathbf{g} \\
\mathbf{a}_{\text{linear,world}} &= \mathbf{R}\,\mathbf{a}_{\text{linear,body}}
\end{align}

The quaternion-to-rotation-matrix:
\begin{equation}
\mathbf{R} = \begin{bmatrix}
1-2(y^2+z^2) & 2(xy-wz) & 2(xz+wy) \\
2(xy+wz) & 1-2(x^2+z^2) & 2(yz-wx) \\
2(xz-wy) & 2(yz+wx) & 1-2(x^2+y^2)
\end{bmatrix}
\end{equation}

% ----------------------------------------------------------------------
\subsection{Trapezoidal Numerical Integration}

Trapezoidal rule (more accurate than forward Euler):
\begin{align}
\mathbf{v}[n] &= \mathbf{v}[n\!-\!1] + \frac{\Delta t}{2}\bigl(\mathbf{a}[n] + \mathbf{a}[n\!-\!1]\bigr) \\
\mathbf{p}[n] &= \mathbf{p}[n\!-\!1] + \frac{\Delta t}{2}\bigl(\mathbf{v}[n] + \mathbf{v}[n\!-\!1]\bigr)
\end{align}
$\mathcal{O}(\Delta t^3)$ local error vs.\ $\mathcal{O}(\Delta t^2)$ for Euler, at 100\,Hz ($\Delta t = 10$\,ms).

% ----------------------------------------------------------------------
\subsection{Zero-Velocity Update (ZUPT)}

Bounds unbounded drift from double integration:
$$\text{If } \|\mathbf{a}_{\text{filtered}}\| < 0.15\text{ m/s}^2 \text{ for 10 consecutive samples} \implies \mathbf{v} \leftarrow \mathbf{0}$$
Filter states are also reset to prevent residual bias accumulation.

% ----------------------------------------------------------------------
\subsection{Spectral Analysis (FFT/PSD)}

\textbf{Hann window} reduces spectral leakage:
$w[n] = 0.5\bigl(1 - \cos(2\pi n/(N-1))\bigr)$, achieving $-31$\,dB sidelobes.

\textbf{PSD} via periodogram:
\begin{equation}
\text{PSD}[k] = \frac{|X[k]|^2}{f_s \cdot \sum w^2[n]}, \qquad
\text{PSD}_{\text{dB}} = 10\log_{10}(\text{PSD})
\end{equation}

Metrics: noise floor (median PSD), peak frequency, SNR (peak $-$ floor), confidence (peak/mean).

% ======================================================================
\section{Software Implementation}

\subsection{Node Summary}

\begin{center}
\begin{tabular}{|l|l|p{5.5cm}|}
\hline
\textbf{Node} & \textbf{Executable} & \textbf{Function} \\
\hline
HTTPReceiverNode     & http\_receiver     & HTTP server (port 5555) for phone JSON data \\
\hline
MadgwickFilterNode   & madgwick\_filter   & AHRS fusion ($\beta\!=\!0.033$, 100\,Hz) \\
\hline
SpectralAnalysisNode & spectral\_analysis & 256-pt Hann-windowed FFT/PSD \\
\hline
InertialNavNode      & inertial\_nav      & Gravity removal, LPF, integration, ZUPT \\
\hline
WSBridgeNode         & ws\_bridge         & WebSocket server (port 8765) \\
\hline
\end{tabular}
\end{center}

\subsection{Web Dashboard}
Connects via WebSocket and displays: 3D orientation (Canvas + Euler angles), acceleration time-series, velocity/displacement plots, FFT spectrum with noise floor/SNR metrics, and pipeline status indicators.

% ======================================================================
\section{Testing}

Comprehensive \texttt{pytest} suite runs without ROS\,2 via mock fixtures in \texttt{conftest.py}:

\begin{center}
\begin{tabular}{|l|p{7cm}|}
\hline
\textbf{Test File} & \textbf{Coverage} \\
\hline
test\_dsp\_filters.py      & Coefficient design, Nyquist check, DC passthrough, noise attenuation \\
\hline
test\_madgwick\_filter.py   & Calibration, quaternion normalization, stationary orientation \\
\hline
test\_inertial\_nav.py      & Gravity removal, trapezoidal integration, ZUPT \\
\hline
test\_spectral\_analysis.py & PSD, windowing, peak detection, SNR \\
\hline
test\_http\_receiver.py     & JSON parsing, CORS headers, error handling \\
\hline
\end{tabular}
\end{center}

% ======================================================================
\section{Results and Discussion}

\subsection{Filter Performance}
Butterworth LPF ($f_c\!=\!5$\,Hz, $f_s\!=\!100$\,Hz): unity DC gain, $-3$\,dB at 5\,Hz, $-40$\,dB/dec rolloff. Effective MEMS noise suppression above 10\,Hz.

\subsection{Orientation}
Madgwick converges within 0.5\,s. Computationally lighter than Kalman with comparable 6-axis accuracy.

\subsection{Navigation}
Trapezoidal integration is smoother than Euler. ZUPT bounds drift during stationary periods. Long-term position drift is inherent without external aiding (GPS, visual odometry).

\subsection{Limitations}
(1)~Drift from double integration of biased sensors. (2)~No magnetometer: yaw drifts freely. (3)~HTTP latency: 1--10\,ms on Wi-Fi.

% ======================================================================
\section{Conclusion}

This project demonstrates a complete real-time inertial navigation pipeline from DSP first principles: Butterworth IIR filtering via bilinear transform, Madgwick AHRS fusion, quaternion-based gravity removal, trapezoidal integration with ZUPT, and Hann-windowed FFT spectral analysis---all within a modular ROS\,2 architecture with live web visualization.

% ======================================================================
\section*{References}
\begin{enumerate}
  \item[{[1]}] A.\,V.\ Oppenheim and A.\,S.\ Willsky, \textit{Signals and Systems}, 2nd ed., Prentice Hall, 1997.
  \item[{[2]}] S.\,O.\,H.\ Madgwick, ``An efficient orientation filter for inertial and inertial/magnetic sensor arrays,'' Univ.\ of Bristol, 2010.
  \item[{[3]}] O.\,J.\ Woodman, ``An introduction to inertial navigation,'' Univ.\ of Cambridge, UCAM-CL-TR-696, 2007.
  \item[{[4]}] G.\ Heinzel et al., ``Spectrum and spectral density estimation by the DFT,'' Max Planck Institute, 2002.
  \item[{[5]}] ROS\,2 Humble Docs, \texttt{https://docs.ros.org/en/humble/}.
\end{enumerate}

% ======================================================================
\appendix
\section{How to Run}
\begin{verbatim}
cd ~/phone_imu_ws
colcon build --symlink-install
source install/setup.bash
ros2 launch phone_imu_bridge imu_bridge.launch.py
# Phone: HTTP Push URL = http://<PC-IP>:5555/data
# Dashboard: open webapp/index.html in browser
\end{verbatim}

\section{Project File Structure}
\begin{verbatim}
phone_imu_bridge/
  phone_imu_bridge/           # Python package
    dsp_filters.py            # Butterworth LPF (first principles)
    http_receiver_node.py     # HTTP server for phone data
    madgwick_filter_node.py   # AHRS sensor fusion
    inertial_nav_node.py      # Dead-reckoning navigation
    spectral_analysis_node.py # FFT/PSD analysis
    ws_bridge_node.py         # ROS-to-WebSocket bridge
  webapp/                     # Web dashboard
    index.html, phone.html
    css/style.css
    js/app.js, plots.js, orientation3d.js, phone.js
  launch/imu_bridge.launch.py
  config/params.yaml
  test/conftest.py, test_*.py # pytest suite
  package.xml, setup.py, README.md
\end{verbatim}

\end{document}
